% ============================================================================
% MODULE 3: OUTPUT ANALYSIS
% ============================================================================
% Duration: 60 minutes (data analysis, visualization)
% Learning Goals: Extract insights from simulation output; plotting; Excel skills
% Prerequisites: Module 2

\chapter{Module 3: Output Analysis}
\label{ch:module3}

\section{Learning Objectives}

By the end of this module, you will:
\begin{itemize}
  \item Locate and export simulation results from the APSIM \apsimterm{DataStore}
  \item Import APSIM output into Excel and Python
  \item Create summary statistics (mean, min, max, standard deviation)
  \item Generate time-series and bar plots from daily output
  \item Interpret crop growth curves (LAI, biomass, grain mass phases)
  \item Calculate key agronomic metrics: yield, NUE, and Harvest Index
\end{itemize}

\section{APSIM Output Structure}

Each row in an APSIM output table represents one simulated day. Columns correspond to the variables you selected in the \apsimterm{Report} component.

\begin{center}
\begin{tabular}{l|r|r|r|r}
\textbf{Date} & \textbf{ESW (mm)} & \textbf{Rainfall (mm)} & \textbf{Drainage (mm)} & \textbf{Evaporation (mm)} \\
\hline
1950-01-01 & 250 & 0 & 0 & 3.5 \\
1950-01-02 & 246 & 0 & 0 & 4.1 \\
1950-01-03 & 257 & 15 & 0 & 2.1 \\
\ldots & \ldots & \ldots & \ldots & \ldots \\
\end{tabular}
\end{center}

\begin{keyconceptbox}
\textbf{Output = daily snapshot.} Every column is a state variable recorded at the end of that day's calculation cycle. If a variable is missing from your output, check the \apsimterm{Report} component's variable list and re-run the simulation.
\end{keyconceptbox}

\subsection{Variable Naming Convention}

APSIM uses dot notation to identify variables:

\begin{verbatim}
[ComponentName].[SubComponent].[Variable]

Example:  [Soil].SoilWater.ESW
Example:  [Wheat].GrainMass
Example:  [Weather].Rainfall
\end{verbatim}

The square brackets identify the component name; the dot separates levels of the hierarchy. These names appear as column headers in exported data.

\section{Exporting Data from APSIM}

\subsection{Step-by-Step: Export to Excel}

\begin{enumerate}
  \item Run your simulation (press \textbf{Run} or F5)
  \item In the \apsimterm{Simulations} tree on the left, click your simulation node
  \item In the right-hand panel, click the \textbf{Results} tab at the bottom of the screen
  \item Right-click anywhere in the results table
  \item Select \textbf{Export} $\to$ \textbf{Excel (.xlsx)}
  \item Save the file to your working folder
\end{enumerate}

% ---- IMAGE PLACEHOLDER ----
\begin{figure}[H]
  \centering
  \fbox{\parbox{0.85\textwidth}{\centering\vspace{1.5cm}
    \textit{[IMAGE: Screenshot of APSIM with the Results tab selected and the right-click Export menu visible.}\\
    \textit{Annotate: (1) Simulation tree, (2) Results tab at bottom, (3) Right-click Export option]}
  \vspace{1.5cm}}}
  \caption{Exporting simulation output from the APSIM DataStore. Select the Results tab, right-click the table, and choose Export.}
  \label{fig:datastore-export}
\end{figure}
% ---- END PLACEHOLDER ----

\begin{tipbox}
\textbf{CSV alternative:} If Excel is not available, select \textbf{CSV (.csv)} instead. CSV files open in any text editor and import cleanly into Python, R, or Jupyter without modification.
\end{tipbox}

\section{Excel Analysis}

\subsection{Opening the Exported File}

\begin{enumerate}
  \item Open the exported \filename{.xlsx} file in Excel
  \item Confirm the sheet name matches your simulation name
  \item Check the first and last rows for date range (should match your Clock start/end dates)
  \item Scan the column headers --- they should match the variable names from your \apsimterm{Report}
\end{enumerate}

\begin{warningbox}
\textbf{Common issue:} The Date column may import as plain text rather than a date.\\
\textit{Fix:} Select the Date column $\to$ Data $\to$ Text to Columns $\to$ select Date (YMD) $\to$ Finish.
\end{warningbox}

\subsection{Summary Statistics with Excel Formulas}

Once your data is loaded, calculate basic statistics for any numeric column. For a column in cells B2:B2558 (seven years of daily data):

\begin{center}
\begin{tabular}{l|l}
\textbf{Statistic} & \textbf{Excel Formula} \\
\hline
Mean & \texttt{=AVERAGE(B2:B2558)} \\
Minimum & \texttt{=MIN(B2:B2558)} \\
Maximum & \texttt{=MAX(B2:B2558)} \\
Std Deviation & \texttt{=STDEV(B2:B2558)} \\
Days with Rainfall $>$ 0 & \texttt{=COUNTIF(B2:B2558,">0")} \\
Days with ESW $>$ 100 mm & \texttt{=COUNTIF(C2:C2558,">100")} \\
\end{tabular}
\end{center}

\begin{tipbox}
\textbf{Tip:} Create a separate summary sheet. Reference raw data in that sheet (e.g., \texttt{=Sheet1!B2:B2558}) so your original data stays unchanged and the summary is easy to update when you switch simulations.
\end{tipbox}

\subsection{Monthly Averages Using a PivotTable}

\begin{enumerate}
  \item Click anywhere in your data table
  \item Select \textbf{Insert} $\to$ \textbf{PivotTable} $\to$ \textbf{New Worksheet}
  \item In the PivotTable field list:
    \begin{itemize}
      \item Drag \textbf{Date} to the \textbf{Rows} area
      \item Right-click any date cell in the table $\to$ \textbf{Group} $\to$ tick \textbf{Months} and \textbf{Years} $\to$ OK
      \item Drag \textbf{ESW} (or your variable of interest) to the \textbf{Values} area
      \item In the Values area, click the drop-down arrow $\to$ \textbf{Value Field Settings} $\to$ choose \textbf{Average}
    \end{itemize}
  \item You now have average ESW for each month-year combination
\end{enumerate}

% ---- IMAGE PLACEHOLDER ----
\begin{figure}[H]
  \centering
  \fbox{\parbox{0.85\textwidth}{\centering\vspace{1.5cm}
    \textit{[IMAGE: Excel screenshot showing a PivotTable with Date (grouped by Month/Year) in rows and ``Average of ESW'' in values.}\\
    \textit{Annotate: (1) Field list panel, (2) Date grouping dialog, (3) Resulting summary table]}
  \vspace{1.5cm}}}
  \caption{Excel PivotTable for monthly average ESW. Group the Date field by Month and Year to produce seasonal summaries without writing formulas.}
  \label{fig:excel-pivot}
\end{figure}
% ---- END PLACEHOLDER ----

\subsection{Visualisation in Excel}

\subsubsection{Time-Series Line Chart (ESW Over 7 Years)}

\begin{enumerate}
  \item Select the Date column, then hold Ctrl and select the ESW column
  \item Select \textbf{Insert} $\to$ \textbf{Charts} $\to$ \textbf{Line} $\to$ \textbf{Line}
  \item Add a title: ``Fallow Water Balance: ESW 1950--1956''
  \item Label axes: X = ``Date'', Y = ``ESW (mm)''
  \item Add a reference line at ESW = 100 mm (sowing threshold):
    \begin{itemize}
      \item In a spare column, enter 100 for every row
      \item Right-click the chart $\to$ \textbf{Select Data} $\to$ \textbf{Add Series}
      \item Select the column of 100s; set series name to ``Sowing threshold''
      \item Format this series as a dashed red line with no markers
    \end{itemize}
\end{enumerate}

\subsubsection{Annual Rainfall Bar Chart}

\begin{enumerate}
  \item In a spare area, create a table with years (1950, 1951, \ldots, 1956) in column A
  \item In column B, calculate annual rainfall using:\\
    \texttt{=SUMPRODUCT((YEAR(\$A\$2:\$A\$2558)=D2)*\$B\$2:\$B\$2558)}\\
    where D2 contains the year value
  \item Select the year and rainfall columns $\to$ Insert $\to$ Column Chart
\end{enumerate}

\begin{tipbox}
\textbf{Conditional formatting:} Highlight sowing-opportunity rows automatically.\\
Select the ESW column $\to$ Home $\to$ Conditional Formatting $\to$ Highlight Cell Rules $\to$ Greater Than $\to$ enter 100 $\to$ choose green fill.
\end{tipbox}

\section{Python Analysis}

\begin{tipbox}
\textbf{Interactive Notebook Available:} A complete Jupyter notebook with all Python analysis code is provided as \texttt{APSIM\_Analysis\_Notebook.ipynb}. Students can:
\begin{itemize}
  \item Open it directly in Google Colab (no installation required)
  \item Upload APSIM CSV exports and run analysis interactively
  \item Ask AI assistants (Gemini, ChatGPT, Claude) for help with code or interpretation
  \item Modify code to suit their specific simulation scenarios
\end{itemize}
Access the notebook from the course LMS or GitHub repository.
\end{tipbox}

\subsection{Reading Data}

\begin{lstlisting}[language=Python, caption=Read APSIM output in pandas]
import pandas as pd
import matplotlib.pyplot as plt

# Read exported CSV
df = pd.read_csv('Exercise_A_Results.csv')

# Convert date column
df['Date'] = pd.to_datetime(df['Date'])

# Examine first few rows
print(df.head())
\end{lstlisting}

\subsection{Summary Statistics}

\begin{lstlisting}[language=Python, caption=Calculate monthly averages]
# Extract year and month
df['YearMonth'] = df['Date'].dt.to_period('M')

# Monthly average ESW
monthly_esw = df.groupby('YearMonth')['ESW'].mean()
print(monthly_esw)
\end{lstlisting}

\subsection{Visualization}

\begin{lstlisting}[language=Python, caption=Plot ESW time series]
plt.figure(figsize=(12, 5))
plt.plot(df['Date'], df['ESW'], linewidth=1)
plt.axhline(y=100, color='r', linestyle='--', label='Sowing threshold')
plt.xlabel('Date')
plt.ylabel('ESW (mm)')
plt.title('Fallow Water Balance: 7-Year ESW Trend')
plt.legend()
plt.grid()
plt.show()
\end{lstlisting}



\section{Agronomic Metrics}
\label{sec:agronomic-metrics}

\subsection{Yield}

APSIM reports \variable{[Wheat].GrainMass} in kg/ha at harvest. Use this value directly. Typical dryland wheat yields in subtropical Queensland: 1500--4500 kg/ha depending on rainfall.

\begin{tipbox}
\textbf{Note:} Some APSIM versions report grain mass in g/m$^2$. To convert to kg/ha: multiply by 10. Always check the Report variable description if yield values seem unexpectedly high or low.
\end{tipbox}

\subsection{Nitrogen Use Efficiency (NUE)}

\[
\text{AEN} = \frac{\text{Yield}_{N \text{ treatment}} - \text{Yield}_{\text{0 N}}}{\text{N rate}}
\]

Typical range: 10--20 kg grain per kg N applied

\subsection{Harvest Index}

\[
\text{HI} = \frac{\text{Grain mass}}{\text{Total aboveground biomass}}
\]

Typical wheat: 0.4--0.5 (40--50\% of biomass is grain)

\section{Interpreting Crop Growth Curves}

When you plot daily crop variables against date, four distinct phases become visible:

\begin{center}
\begin{tabular}{l|l|l}
\textbf{Phase} & \textbf{Zadok stages} & \textbf{What you see in output} \\
\hline
Establishment & 0--19 & LAI rises from 0; Biomass near zero \\
Vegetative & 20--49 & LAI increases rapidly; Biomass growing \\
Reproductive & 50--69 & LAI at maximum; Biomass growth peaks \\
Grain fill & 70--89 & LAI declining; GrainMass increases sharply \\
Maturity & 90+ & LAI = 0; GrainMass plateaus (harvest triggered) \\
\end{tabular}
\end{center}

% ---- IMAGE PLACEHOLDER ----
\begin{figure}[H]
  \centering
  \fbox{\parbox{0.85\textwidth}{\centering\vspace{1.5cm}
    \textit{[IMAGE: Dual-axis plot of APSIM output: Biomass (kg/ha, left axis) and LAI (m$^2$/m$^2$, right axis) vs Date.}\\
    \textit{Annotate key phases with vertical dashed lines: Emergence, Anthesis, Grain fill start, Maturity]}
  \vspace{1.5cm}}}
  \caption{Typical wheat crop growth curve from APSIM. Biomass grows continuously until maturity; LAI peaks around anthesis then declines as leaves senesce.}
  \label{fig:growth-curve}
\end{figure}
% ---- END PLACEHOLDER ----

\begin{keyconceptbox}
\textbf{Use the Zadok stage column as a guide.} In your exported data, find rows where \variable{[Wheat].Stage} $\approx$ 60 (anthesis) and \variable{[Wheat].Stage} $\approx$ 70 (grain fill start). Check that LAI is near its peak at anthesis and that GrainMass is essentially zero until grain fill begins. If either is missing, re-check your Manager harvest rules and Report variables.
\end{keyconceptbox}

\section{Common Analysis Mistakes}

\begin{warningbox}
\textbf{Mistake 1:} Comparing simulations with different date ranges

\textit{Solution:} Normalize to common period or use per-season metrics (e.g., cumulative rainfall in growing season, not absolute date ranges)
\end{warningbox}

\begin{warningbox}
\textbf{Mistake 2: Not accounting for missing values before sowing}

APSIM reports \texttt{NaN} or blank for crop variables (e.g., \variable{[Wheat].GrainMass}) during the fallow period before the crop is sown. Averaging a column that includes these blanks will underestimate the true mean.

\textit{Fix:} In Excel, use \texttt{AVERAGEIF} to exclude blanks. In Python, use \texttt{df['col'].dropna().mean()}.
\end{warningbox}

\begin{warningbox}
\textbf{Mistake 3: Interpreting daily output as seasonal totals}

The \variable{Rainfall} column is the rainfall on that day (mm/day), not cumulative. Similarly, ESW is the soil water on that day, not total water added.

\textit{Fix:} To get seasonal totals, use \texttt{SUM} (or \texttt{groupby().sum()} in Python) on the relevant date range.
\end{warningbox}

\section{Summary}

\begin{keyconceptbox}
\textbf{Module 3 key points:}
\begin{enumerate}
  \item APSIM output = one row per simulated day; columns = \apsimterm{Report} variables
  \item Export from \apsimterm{DataStore}: Results tab $\to$ right-click $\to$ Export $\to$ Excel or CSV
  \item Excel: verify date format $\to$ PivotTable for monthly averages $\to$ line chart for time series
  \item Python: \texttt{pd.read\_csv()} $\to$ \texttt{groupby()} $\to$ \texttt{matplotlib}
  \item Core metrics: Yield (kg/ha), AEN (kg grain / kg N), Harvest Index (0.4--0.5 for wheat)
  \item Plot crop growth curves and use Zadok stage to verify all phases are present
\end{enumerate}
\end{keyconceptbox}

\section{What's Next?}

Module 4 explores nitrogen fertilisation and N response curves.

% ============================================================================
% END MODULE 3
% ============================================================================
